\chapter{Prerequisites in mathlib}
\label{ch:prereq}

Everything in this chapter exists in mathlib as of September 2026. The nodes are recorded so
that the graph shows what the project rests on and so that API changes upstream are
visible as broken edges.

\begin{definition}[Cochain complexes of modules]
\label{def:cochain-complex}
\lean{CochainComplex}
\mathlibok
\leanok
For a commutative ring $k$, $\Ch(k)$ denotes the category of $\mathbb{Z}$-indexed cochain
complexes of $k$-modules, \texttt{CochainComplex (ModuleCat k) }$\mathbb{Z}$. Its objects are
Keller's dg $k$-modules.
\end{definition}

\begin{theorem}[Monoidal structure on complexes]
\label{thm:complexes-monoidal}
\lean{HomologicalComplex.monoidalCategory}
\mathlibok
\leanok
\uses{def:cochain-complex}
$\Ch(k)$ is a monoidal category with the tensor product of complexes
$(V \otimes W)^n = \bigoplus_{p+q=n} V^p \otimes W^q$ and differential
$d_V \otimes 1 + 1 \otimes d_W$ with Koszul signs, and unit the complex $k$ in degree $0$.
\end{theorem}

\begin{definition}[Enriched categories]
\label{def:enriched-category}
\lean{CategoryTheory.EnrichedCategory}
\mathlibok
\leanok
For a monoidal category $V$, a $V$-enriched category $\cC$ consists of a type of objects,
hom-objects $\cC(X,Y) \in V$, identities $\mathbf{1} \to \cC(X,X)$, and compositions
$\cC(X,Y) \otimes \cC(Y,Z) \to \cC(X,Z)$ satisfying unit and associativity axioms.
\end{definition}

\begin{definition}[Enriched functors]
\label{def:enriched-functor}
\lean{CategoryTheory.EnrichedFunctor}
\mathlibok
\leanok
\uses{def:enriched-category}
A $V$-enriched functor is a map on objects together with morphisms
$\cC(X,Y) \to \cD(FX,FY)$ in $V$ compatible with identities and composition.
\end{definition}

\begin{definition}[Lax monoidal functors]
\label{def:lax-monoidal}
\lean{CategoryTheory.Functor.LaxMonoidal}
\mathlibok
\leanok
A functor $F : V \to W$ between monoidal categories is lax monoidal if it carries
structure maps $\mathbf{1} \to F\mathbf{1}$ and $FX \otimes FY \to F(X \otimes Y)$
satisfying the usual coherence conditions.
\end{definition}

\begin{theorem}[Transport of enrichment]
\label{thm:transport-enrichment}
\lean{CategoryTheory.TransportEnrichment}
\mathlibok
\leanok
\uses{def:enriched-category, def:lax-monoidal}
A lax monoidal functor $F : V \to W$ turns a $V$-enriched category into a $W$-enriched
category with the same objects and hom-objects $F(\cC(X,Y))$.
\end{theorem}

\begin{theorem}[Forgetting enrichment]
\label{thm:forget-enrichment}
\lean{CategoryTheory.ForgetEnrichment}
\mathlibok
\leanok
\uses{def:enriched-category}
A $V$-enriched category has an underlying ordinary category whose morphisms $X \to Y$ are
the morphisms $\mathbf{1} \to \cC(X,Y)$ in $V$.
\end{theorem}

\begin{definition}[The Hom complex]
\label{def:hom-complex}
\lean{CochainComplex.HomComplex}
\mathlibok
\leanok
\uses{def:cochain-complex}
For cochain complexes $L, M$ in a preadditive category, $\Hom(L,M)$ is the complex whose
degree-$n$ part consists of the families $(f^p : L^p \to M^{p+n})_p$ and whose
differential is $d(f) = d_M \circ f - (-1)^n f \circ d_L$.
\end{definition}

\begin{definition}[Homotopy category of complexes]
\label{def:homotopy-category}
\lean{HomotopyCategory}
\mathlibok
\leanok
\uses{def:cochain-complex}
$\Hcat(B)$ is the quotient of $\Ch(B)$ by the ideal of morphisms homotopic to zero.
\end{definition}

\begin{definition}[Localization of categories]
\label{def:localization}
\lean{CategoryTheory.Functor.IsLocalization}
\mathlibok
\leanok
Given a class of morphisms $W$ in a category $\cC$, a functor $L : \cC \to \cD$ is a
localization of $\cC$ at $W$ if it inverts $W$ and is universal among such functors.
mathlib provides both the universal property and a construction, together with
criteria (calculus of fractions) for the localized category to have small hom-sets.
\end{definition}

\begin{definition}[Derived category of an abelian category]
\label{def:mathlib-derived-category}
\lean{DerivedCategory}
\mathlibok
\leanok
\uses{def:cochain-complex, def:homotopy-category, def:localization}
For an abelian category $\cB$, mathlib's $\Dcat(\cB)$ is the localization of
$\Ch(\cB)$ at quasi-isomorphisms, equivalently the localization of $\Hcat(\cB)$; it
carries a triangulated structure. In Chapter~\ref{ch:derived} this is the target of the
comparison theorem when $\cB = \Mod B$.
\end{definition}
